\Needspace{17\baselineskip}
\Article{21}

Second, let $n$ be the square or a higher power of an odd prime number $p$, and put $n=p^{2\varkappa}q$, with $q$ either $=1$ or $=p$.\quad Here we should first note that, if $\lambda$ is any integer not divisible by $p^\varkappa$, we have

\SourceLeaf{43}

\[
\begin{aligned}
&r^{\lambda\lambda}+r^{(\lambda+p^\varkappa q)^2}+r^{(\lambda+2p^\varkappa q)^2}+r^{(\lambda+3p^\varkappa q)^2}\\&\qquad+\text{etc.}+r^{(\lambda+n-p^\varkappa q)^{3}}\\
&\quad=r^{\lambda\lambda}\{1+r^{2\lambda p^\varkappa q}+r^{4\lambda p^\varkappa q}\\&\qquad+r^{6\lambda p^\varkappa q}+\text{etc.}+r^{2\lambda(n-p^\varkappa q)}\}\\&\quad=\frac{r^{\lambda\lambda}(1-r^{2\lambda n})}{1-r^{2\lambda p^\varkappa q}}=0
\end{aligned}
\]

It is then easy to see that

\[
\begin{aligned}
W=1+r^{p^{2\varkappa}}+r^{4p^{2\varkappa}}+r^{9p^{2\varkappa}}+\text{etc.}+r^{(n-p^\varkappa)^2}
\end{aligned}
\]

\Needspace{6\baselineskip}
For the remaining terms of the series

\[
\begin{aligned}
1+r+r^4+r^9+\text{etc.}+r^{(n-1)^2}
\end{aligned}
\]

can be divided into $(p^\varkappa-1)q$ partial series, each with $p^\varkappa$ terms, whose sums vanish under the transformation just given.

It follows that in the case where $q=1$, or where $n$ is a power of a prime number with an even exponent, we have

\[
\begin{aligned}
W=p^\varkappa=+\surd n,\;\text{and so}\;T=+\surd n,\quad U=0
\end{aligned}
\]

In the other case, where $q=p$, or where $n$ is a power of a prime number with an odd exponent, we shall put $r^{p^{2\varkappa}}=\rho$. Then $\rho$ is a proper root of the equation $x^p-1=0$, namely $\rho=\cos\frac{k}{p}360^0+i\sin\frac{k}{p}360^0$, and we obtain

\[
\begin{aligned}
W&=1+\rho+\rho^4+\rho^9+\text{etc.}+\rho^{(p^{\varkappa+1}-1)^2}\\&=p^\varkappa(1+\rho+\rho^4+\rho^9+\text{etc.}+\rho^{(p-1)^2})
\end{aligned}
\]

But the sum of the series $1+\rho+\rho^4+\rho^9+\text{etc.}+\rho^{(p-1)^2}$ is determined by the preceding article. It follows directly that

\[
\begin{aligned}
W&=\pm\surd n=T,\quad\text{if }p\text{ is of the form }4\mu+1\\
W&=\pm i\surd n=iU,\quad\text{if }p\text{ is of the form }4\mu+3
\end{aligned}
\]

with the positive or negative sign, according as $k$ is a residue or a nonresidue of $p$.

\Needspace{12\baselineskip}
\Article{22}

What we set out in Articles 20 and 21 also readily gives the following proposition, which will be of considerable use to us below.\quad Put

\[
\begin{aligned}
W'=1+r^h+r^{4h}+r^{9h}+\text{etc.}+r^{h(n-1)^2}
\end{aligned}
\]

\SourceLeaf{44}

with $h$ denoting any integer not divisible by $p$. Then, in the case where $n=p$, or where $n$ is a power of $p$ with an odd exponent, we have

\[
\begin{aligned}
W'&=W,\;\text{if }h\text{ is a quadratic residue of }p\\
W'&=-W,\;\text{if }h\text{ is a quadratic nonresidue of }p
\end{aligned}
\]

For $W'$ evidently arises from $W$ if, for $k$, we substitute $kh$. In the first case, $k$ and $kh$ are alike, and in the second they are unlike, with respect to being residues or nonresidues of $p$.

But in the case where $n$ is a power of $p$ with an even exponent, we evidently have $W'=+\surd n$, and therefore always $W'=W$.

\Needspace{13\baselineskip}
\Article{23}

In Articles 20, 21 and 22 we considered odd prime numbers and powers of such numbers. The case therefore remains where $n$ is a power of two.

For $n=2$, we evidently have $W=1+r=0$.

For $n=4$, we obtain $W=1+r+r^4+r^9=2+2r$. Thus $W=2+2i$ whenever $k$ is of the form $4\mu+1$, and $W=2-2i$ whenever $k$ is of the form $4\mu+3$.

For $n=8$, we have $W=1+r+r^4+r^9+r^{16}+r^{25}+r^{36}+r^{49}=2+4r+2r^4$ $=4r$.\quad Hence

\[
\begin{aligned}
W&=(1+i)\surd8,\quad\text{whenever }k\text{ is of the form }8\mu+1\\
W&=(-1+i)\surd8,\quad\text{whenever }k\text{ is of the form }8\mu+3\\
W&=(-1-i)\surd8,\quad\text{whenever }k\text{ is of the form }8\mu+5\\
W&=(1-i)\surd8,\quad\text{whenever }k\text{ is of the form }8\mu+7
\end{aligned}
\]

If $n$ is a higher power of two, put $n=2^{2\varkappa}q$, with $q$ either $=1$ or $=2$, and $\varkappa$ greater than $1$.\quad Here we must first observe that, if $\lambda$ is any integer not divisible by $2^{\varkappa-1}$, we have

\[
\begin{aligned}
&r^{\lambda\lambda}+r^{(\lambda+2^\varkappa q)^2}+r^{(\lambda+2\pdot2^\varkappa q)^2}+r^{(\lambda+3\pdot2^\varkappa q)^2}\\&\qquad+\text{etc.}+r^{(\lambda+n-2^\varkappa q)^2}\\
&\quad=r^{\lambda\lambda}\{1+r^{2^{\varkappa+1}\lambda q}+r^{2\pdot2^{\varkappa+1}\lambda q}\\&\qquad+r^{3\pdot2^{\varkappa+1}\lambda q}+\text{etc.}+r^{(2n-2^{\varkappa+1}q)\lambda}\}\\&\quad=\frac{r^{\lambda\lambda}(1-r^{2\lambda n})}{1-r^{2^{\varkappa+1}\lambda q}}=0
\end{aligned}
\]

It is then easy to see that

\[
\begin{aligned}
W=1+r^{2^{2\varkappa-2}}+r^{4\pdot2^{2\varkappa-2}}+r^{9\pdot2^{2\varkappa-2}}+\text{etc.}+r^{(n-2^{\varkappa-1})^2}
\end{aligned}
\]

\SourceLeaf{45}

Put $r^{2^{2\varkappa-2}}=\rho$. Then $\rho$ is a root of the equation $x^{4q}-1=0$, namely $\rho=\cos\frac{k}{4q}360^0+i\sin\frac{k}{4q}360^0$; we then obtain

\[
\begin{aligned}
W&=1+\rho+\rho^4+\rho^9+\text{etc.}+\rho^{(2^{\varkappa+1}q-1)^2}\\
&=2^{\varkappa-1}(1+\rho+\rho^4+\rho^9+\text{etc.}+\rho^{(4q-1)^2})
\end{aligned}
\]

\Needspace{14\baselineskip}
But the sum of the series $1+\rho+\rho^4+\rho^9+\text{etc.}+\rho^{(4q-1)^2}$ is determined by what we explained for the cases $n=4$, $n=8$. From this we conclude

that in the case where $q=1$, or where $n$ is a power of the number $4$, we have

\[
\begin{aligned}
W&=(1+i)2^\varkappa=(1+i)\surd n,\quad\text{if }k\text{ is of the form }4\mu+1\\
W&=(1-i)2^\varkappa=(1-i)\surd n,\quad\text{if }k\text{ is of the form }4\mu+3
\end{aligned}
\]

these are exactly the formulas given for $n=4$;

but in the case where $q=2$, or where $n$ is a power of two with an odd exponent greater than $3$, we have

\[
\begin{aligned}
W&=(1+i)2^\varkappa\surd2=(1+i)\surd n,\quad\text{if }k\text{ is of the form }8\mu+1\\
W&=(-1+i)2^\varkappa\surd2=(-1+i)\surd n,\quad\text{if }k\text{ is of the form }8\mu+3\\
W&=(-1-i)2^\varkappa\surd2=(-1-i)\surd n,\quad\text{if }k\text{ is of the form }8\mu+5\\
W&=(1-i)2^\varkappa\surd2=(1-i)\surd n,\quad\text{if }k\text{ is of the form }8\mu+7
\end{aligned}
\]

these too agree exactly with those we gave for $n=8$.

\Needspace{12\baselineskip}
\Article{24}

Here too it will be worthwhile to determine the ratio of the sum of the series

\[
\begin{aligned}
W'=1+r^h+r^{4h}+r^{9h}+\text{etc.}+r^{h(n-1)^2}
\end{aligned}
\]

to $W$, where $h$ denotes any odd integer.\quad Since $W'$ arises from $W$ by changing $k$ into $kh$, the value of $W'$ depends on the form of the number $kh$ in the same way as $W$ depends on the form of $k$.\quad Put $\frac{W'}{W}=l$. It is clear that

I.\quad in the case where $n=4$, or it is a higher power of two with an even exponent, we have

\[
\begin{aligned}
l&=1,\quad\text{if }h\text{ is of the form }4\mu+1\\
l&=-i,\quad\text{if }h\text{ is of the form }4\mu+3,\;\text{and }k\text{ is of the form }4\mu+1\\
l&=+i,\quad\text{if }h\text{ is of the form }4\mu+3,\;\text{and }k\text{ is of the same form.}
\end{aligned}
\]

\SourceLeaf{46}

\Needspace{15\baselineskip}
II.\quad in the case where $n=8$, or it is a higher power of two with an odd exponent, we have

\[
\begin{aligned}
l&=1,\ \text{if}\  h\text{ is of the form }8\mu+1,\\
l&=-1,\ \text{if}\  h\text{ is of the form }8\mu+5,\\
l&=+i,\ \text{if either }h\text{ is of the form }8\mu+3,\;\text{and }k\text{ is of the form }4\mu+1,\\
&\phantom{=+i,}\text{or }h\text{ is of the form }8\mu+7,\;\text{and }k\text{ is of the form }4\mu+3,\\
l&=-i,\ \text{if either }h\text{ is of the form }8\mu+3,\;\text{and }k\text{ is of the form }4\mu+3,\\
&\phantom{=+i,}\text{or }h\text{ is of the form }8\mu+7,\;\text{and }k\text{ is of the form }4\mu+1.
\end{aligned}
\]

The preceding articles have completely determined the sum $W$ for the cases where $n$ is a prime number or a power of a prime number. It therefore remains to complete the cases where $n$ is composed of several prime numbers; the following theorem will prepare the way for us:
